\section{Keeping Advent}

\subsection{Watch}\label{sec:watch}

The first word of Advent is a command. In every year of the cycle the Gospel of
the First Sunday tells the Church to watch. In Year B it is the end of the Lord's last discourse in St Mark:

\begin{quote}
Take ye heed, watch and pray. For ye know not when the time is. Even as a man
who, going into a far country, left his house and gave authority to his servants
over every work and commanded the porter to watch. Watch ye therefore (for you
know not when the lord of the house cometh, at even, or at midnight, or at the
cock crowing, or in the morning): Lest coming on a sudden, he find you sleeping.
And what I say to you, I say to all: Watch.\par
{\smallskip\noindent\small ---Mark 13:33--37\par}
\end{quote}

\noindent In Year A the Church hears the same charge from St Matthew: ``Watch ye
therefore, because you know not what hour your Lord will come.'' In Year C she
hears it from St Luke: ``Watch ye, therefore, praying at all times, that you may
be accounted worthy to escape all these things that are to come and to stand
before the Son of man.''

St Mark's four watches of the night, evening, midnight, cockcrow and morning,
were read by the Greek commentator Theophylact as the ages of a life, as the
Catena Aurea gives him:

\begin{quote}
Let us therefore attend to His words; for the end comes at even, when a man dies
in old age; a midnight, when he dies in the midst of his youth; and at cockcrow,
when our reason is perfect within us; for when a child begins to live according
to his reason, then the cock cries loud within him, rousing him from the sleep
of sense; but the age of childhood is the morning. Now all these ages must look
out for the end\par
{\smallskip\noindent\small ---Theophylact, in the Catena Aurea on Mark 13:35\par}
\end{quote}

Why is the hour hidden at all? St John Chrysostom, preaching on St Matthew,
found the reason in the Lord's care for his servants:

\begin{quote}
For this intent He tells them not, in order that they may watch, that they may
be always ready; therefore He saith, When ye look not for it, then He will come,
desiring that they should be anxiously waiting, and continually in virtuous
action. \ldots{} In order therefore that they may strive, not at that hour only,
therefore He tells them not either the common hour, or the hour of each,
desiring them to be ever looking for this, that they may be always striving.
Wherefore He made the end of each man's life also uncertain.\par
{\smallskip\noindent\small ---St John Chrysostom, Homily 77 on St Matthew\par}
\end{quote}

\noindent St Augustine put the question himself, how the Lord's words could be
meant for people who would die long before his return, and gave a like answer:

\begin{quote}
For that day comes to each man when his day comes for departing from this life
such as he is to be, when judged in that day, and for this reason every
Christian ought to watch, lest the Advent of the Lord find him unprepared; but
that day shall find him unprepared, whom the last day of his life shall find
unprepared.\par
{\smallskip\noindent\small ---St Augustine, Letter 199, in the Catena Aurea on Mark 13:37\par}
\end{quote}

What is it to watch? Pope St Gregory the Great, preaching on the servants who
wait for their lord's return from the wedding, said it three times over:
\latin{Vigilat qui ad aspectum veri luminis mentis oculos apertos tenet, vigilat
qui servat operando quod credit, vigilat qui a se torporis et negligentiæ
tenebras repellit}. He watches who keeps the eyes of his mind open to the sight
of the true light; he watches who keeps by his deeds what he believes; he
watches who drives from himself the darkness of sloth and neglect. And Gregory
went straight to St Paul's words, which the Church reads on the First Sunday in
Year A: \latin{Hora est jam nos de somno surgere}, it is now the hour to rise
from sleep. In the Douay the verse runs whole:

\begin{quote}
And that, knowing the season, that it is now the hour for us to rise from sleep.
For now our salvation is nearer than when we believed.\par
{\smallskip\noindent\small ---Romans 13:11\par}
\end{quote}

In another homily, on the wise and foolish virgins, Gregory gave the hard edge
of the same truth. God, who receives the penitent after sin, has promised
pardon to the one who repents: \latin{Sed qui pœnitenti veniam spopondit,
peccanti diem crastinum non promisit}. But he who promised pardon to the
penitent did not promise tomorrow to the sinner. \latin{Semper ergo extremum diem
debemus metuere, quem nunquam possumus prævidere}: we must always fear the last
day, because we can never foresee it.

The watch is kept the whole season through. On the Second and Third Sundays
John the Baptist cries out for the way to be made ready, and on the Fourth the
Gospel stands at the threshold of the Lord's birth.
